% Executive Summary — self-contained single page (excluded from D6 page count)

A ground search team covers roughly 0.1--0.5\,km$^{2}$/hr; a single drone at 30\,m altitude covers 0.6\,km$^{2}$/hr---an order-of-magnitude improvement that can compress the \emph{golden hour} search window from hours to minutes. This report presents the design, implementation, and verification of a fully autonomous SAR drone that performs onboard casualty detection using commodity hardware: a Raspberry Pi~5, an RGB camera, and a YOLOv8n neural network---no GPU, no ground-station video link, no operator in the detection loop.

\paragraph{Mission.}
The drone takes off, searches a polygon-defined area via a lawnmower pattern, detects a casualty surrogate with onboard computer vision, reports the target's GPS coordinates and imagery to a ground operator for confirmation, lands 5--10\,m from the casualty to enable first-aid-kit deployment, and returns to the take-off location. A five-layer automatic geofence (waypoint filtering, speed ramp, repulsive field, hard boundary cutoff, and firmware backup) enforces the SSSI no-fly zone, 50\,m altitude ceiling, and flight area boundaries throughout.

\paragraph{Architecture.}
The platform is a Hexsoon EDU-450 hexacopter with a Cube Orange+ autopilot running ArduPilot. All onboard intelligence runs on the Pi~5: a YOLOv8n detector exported to TensorFlow Lite, a lawnmower path planner, GPS-based target localisation with inverse-variance-weighted spatial clustering, and a 20-state finite state machine with 32~transitions orchestrating the complete mission. A browser-based ground station streams live MJPEG video and provides per-target classification controls over Wi-Fi. Four MCDA trade studies informed the selection of detection model, search pattern, companion computer, and ground-station architecture.

Three architectural choices define the system: (i)~a \emph{modular single-codebase} design that runs identical code across laptop simulation, bench hardware, and the Pi without modification; (ii)~a \emph{dual-backend vision pipeline} that automatically selects Ultralytics on the laptop and TFLite on the Pi; and (iii)~a \emph{simulation-first methodology} with a five-tier progressive testing framework---SITL simulation, hardware component checks, bench integration, passive field observation, and autonomous flight---supported by 74~test scripts (including 127~unit tests) across eight categories.

\paragraph{Key results.}
\begin{itemize}[nosep,leftmargin=1.5em]
  \item \textbf{Detection:} mAP$_{50}$ = \textbf{0.995} (retrained on 366 mixed synthetic/real images at native $1456 \times 1088$ resolution).
  \item \textbf{Inference speed:} \textbf{206\,ms} per frame (\textbf{4.8\,FPS}) on the Pi~5 CPU via XNNPACK; lens distortion correction adds only 1.5\,ms.
  \item \textbf{Localisation:} GPS target estimation with CEP$_{50}$ = \textbf{2.3\,m}, confirmed via DJI flight-video replay over the test site (15--50\,m altitude range).
  \item \textbf{Coverage:} Lawnmower pattern covers the survey polygon in $\sim$1.1\,minutes (18~waypoints), verified via dry-run to avoid SSSI encroachment.
  \item \textbf{Requirements:} All twelve brief requirements (R01--R12) addressed and verified through simulation, bench testing, and field-day calibration.
\end{itemize}

\paragraph{Field integration.}
Full hardware integration---Pi, Cube, camera, and GPS---was validated on the assembled airframe during the scheduled field day at Fenswood Farm. On-site bench testing produced lens calibration data (RMS\,=\,0.399), FOV characterisation (focal length 5.46\,mm), and inference benchmarks confirming 100\,\% detection at 0.966 mean confidence across 50~consecutive frames. High winds prevented powered flight on the day; consequently, the complete mission pipeline---arming, search, detection, operator verification, and landing---was exercised end-to-end in SITL simulation, and DJI flight video recorded over the test site served as a real-world proxy for altitude-dependent detection analysis.

\paragraph{Status.}
The system is \textbf{flight-ready}: the full mission executes autonomously in simulation, all hardware subsystems are individually validated on the airframe, and calibration data from the field day is integrated into the deployed pipeline. The single remaining step is powered autonomous flight under favourable weather conditions. The progressive testing methodology, operator-in-the-loop verification, and layered geofence architecture provide quantified confidence in each subsystem for this remaining step.
